% SHARD-ID: AQM-73-QSA-FUSION-CATEGORY-ENDPOINT
% SHARD-TITLE: Fusion-Category Endpoint
% SHARD-SUMMARY: Records fusion categories as a future quantum-system association family without developing the technical construction.
% SHARD-SUMMARY: Lists the source and convention decisions needed before using hom spaces, simple objects, tensor products, F/R data, and lattice-model input.
% SHARD-SUMMARY: Explicitly defers Levin-Wen/string-net claims, fusion-rule computations, Hamiltonians, projectors, and quantum-double or Drinfeld-center statements.
% SHARD-KEYWORDS: quantum-system association, fusion category, gap shard, Levin-Wen, string-net, F-symbol, R-symbol, Drinfeld center

\section{Fusion-Category Endpoint}
\label{sec:qsa-fusion-category-endpoint}

\subsection{Status and Local Anchors}

This is QSA Step 09 from
\path{docs/quantum_system_associations/PLAN.md}.  It is a gap shard, not a
technical fusion-category shard.  It adds no source-local definition of a
fusion category, no example, no fusion-rule calculation, no \(F\)-symbol or
\(R\)-symbol equation, no Levin-Wen/string-net Hamiltonian, and no
comparison theorem with earlier QSA workflows.

The local anchors are the source queue in AQM-04, which still asks for
Levin-Wen string-net models and categorical input data beyond the toric
stabilizer/CSS layer; the QSA meta-plan AQM-64; the vocabulary and evidence
rules of AQM-65 and \texttt{CONVENTIONS.md} convention \texttt{(ap)}; the
finite AND/OR deferral in AQM-70 and convention \texttt{(at)}; and the
boson/fermion gap statement in AQM-72.  Convention \texttt{(h)} explicitly
leaves Levin-Wen fusion-category gauge, pivotal/spherical structure, and
\(F/R\)-data unfixed.

The relevant local source manifest is
\path{references/toric_code/SOURCES.md}.  It currently registers sources for
the toric-code stabilizer and CSS boundary-map conventions used earlier in
the report.  This shard does not treat that manifest as a source for
fusion-category, Levin-Wen, or string-net technical data.

\subsection{Intended Role in the Catalogue}

The intended role of this endpoint is to reserve a future association family
for categorical kinematical data.  Earlier QSA shards cover formal carriers,
tensor-site coexistence, direct-sum alternatives, finite AND/OR carrier
syntax, finite symplectic Weyl labels, and reviewed CCR/CAR layers.  Those
shards also state that para, statistics, Fock, and fusion-category
enhancements are separate choices.

The future fusion-category family is therefore a planned catalogue slot for
workflows in which the input is not merely a set of sites or finite-dimensional
carriers.  It must include a registered categorical source, the category-side
data selected from that source, any lattice or local-constraint data used to
build a kinematical carrier, and the exact evidence status of each output.
Until those items exist locally, the endpoint remains a blocked/gap row in the
QSA catalogue.

\subsection{Required Source and Convention Decisions}

Before any technical statement is added under this endpoint, the following
decisions must be sourced or recorded as conventions.

\begin{center}
\small
\begin{tabular}{p{0.26\linewidth}p{0.64\linewidth}}
\toprule
Decision slot & Required local record before use \\
\midrule
Hom spaces &
The source and notation for hom spaces, their ground field or coefficient
setting, finite-dimensionality assumptions if any, basis choices if any, and
the meaning of adjoints or duals if those are used. \\
\addlinespace
Simple objects &
The source's convention for simple objects, the chosen label set, ordering of
labels, unit label, dual labels, and whether representatives or isomorphism
classes are being used in local tables. \\
\addlinespace
Tensor products and fusion rules &
The tensor-product convention, unit convention, decomposition convention, and
the exact source or checked computation behind any fusion-rule table. \\
\addlinespace
Associator and \(F\)-symbol data &
The associator convention, \(F\)-symbol notation, gauge or basis convention,
normalization choices, and the source or local check for any equation using
that data. \\
\addlinespace
Pivotal or spherical structure &
Whether pivotal, spherical, trace, dimension, or orientation-sensitive data
are required, and which source/convention fixes them. \\
\addlinespace
Braiding and \(R\)-symbols &
Whether braiding is part of the selected workflow.  If it is, the source,
gauge, label order, and \(R\)-symbol convention must be recorded before any
braiding-dependent statement appears. \\
\addlinespace
Lattice, Levin-Wen, or string-net data &
The lattice or cellulation convention, orientation and boundary convention,
where labels live, what local data are attached to vertices/edges/faces, and
which source fixes the model input. \\
\addlinespace
Hilbert carrier and local constraints &
The carrier basis convention, inner-product convention, local-constraint
convention, and evidence that the resulting kinematical subspace or algebra is
the intended output. \\
\bottomrule
\end{tabular}
\end{center}

These are requirement slots, not imported definitions.  A later shard may
split them across several conventions if the selected sources use incompatible
normalizations.

\subsection{Explicit Deferrals}

The following are not claimed anywhere in this shard.

\begin{itemize}
\item No Levin-Wen or string-net technical construction is developed.
\item No fusion-rule table is computed or imported.
\item No \(F\)-symbol equation, \(R\)-symbol equation, gauge transformation,
or coherence check is stated.
\item No Hamiltonian, local projector, commuting-projector result, ground
space, excitation, or topological order claim is stated.
\item No quantum-double, Drinfeld-center, or category-of-excitations claim is
stated.
\item No comparison is made between a future fusion-category workflow and the
AQM-68 tensor-site workflow, the AQM-70 finite AND/OR grammar, or the AQM-72
CCR/CAR layers.
\end{itemize}

Each item remains blocked until a local source manifest, convention entry, and
either a faithful local derivation or reproducible check support it.

\subsection{Unblock Checklist}

To turn this endpoint into a technical association family, a later session
should first acquire and register the relevant fusion-category and
Levin-Wen/string-net sources, then add the convention entries required above.
Only after that should it introduce examples, tables, equations, Hamiltonian
terms, projectors, or comparison maps.  The first technical shard should be
able to point every displayed datum to a local source locator, checked local
derivation, or run artifact.

\subsection{Conclusion}

QSA Step 09 closes as a gap record.  Fusion categories are kept in the
association catalogue as a future endpoint, but the report currently has only
the source queue and convention warnings needed to say what is missing.  The
technical Levin-Wen/string-net and categorical claims are deliberately
deferred.
